\section{The Source-Count Distribution}
\label{sec:source_count}

The analysis of the unresolved gamma-ray sky begins with a single observable: the differential source-count distribution, $dN/dS$.
This function describes how many point sources exist at each flux level and provides the most direct statistical characterization of the populations that contribute to the unresolved gamma-ray background.
Its reconstruction below the telescope detection threshold --- where individual sources can no longer be identified --- is the central inference problem addressed in this chapter.

This section introduces the $dN/dS$ and the methodological context necessary to understand the approach developed in Section~\ref{sec:sbi_cnn}. We begin with the formal definition and its connection to the underlying luminosity function (Section~\ref{sec:6.2.1}), then review the 1-point photon-count distribution, the principal analytical method used to extract $dN/dS$ from sub-threshold data (Section~\ref{sec:6.2.2}). We close by discussing how measurements of the total source-count distribution, combined with population-specific estimates, constrain the composition of the unresolved gamma-ray sky (Section~\ref{sec:6.2.3}).

\subsection{Definition and Physical Meaning}
\label{sec:6.2.1}

The differential source-count distribution, $dN/dS$, gives the number of sources per unit solid angle per unit flux at a given integral photon flux $S$ in a specified energy band.
It is the fundamental observable that describes how point sources are distributed as a function of their apparent brightness, and it connects directly to the physics of the emitting populations.

For sources bright enough to be individually cataloged, $dN/dS$ is measured by simply counting objects in the catalog as a function of their flux (see also Section~\ref{sec:astro_sky}). The 4FGL-DR3 catalog \cite{2022ApJS..260...53A}, based on 12 years of Fermi-LAT data, contains 6658 resolved sources and provides a direct measurement of the source-count distribution in the resolved regime \cite{Amerio:2023uet}.
Above the flux threshold where the catalog detection efficiency reaches 100\%, this measurement is straightforward.
Below this threshold, the detection efficiency degrades, and direct counting must be supplemented by Monte Carlo corrections to reconstruct the true underlying distribution.
Further below --- in the unresolved regime, where individual sources are too faint to be detected at all --- the cumulative emission of these sources forms the unresolved gamma-ray background (UGRB), and the $dN/dS$ must be inferred from the statistical properties of the photon-count map.
This reconstruction is the central challenge addressed by this chapter.

The source-count distribution is the observational projection of a more fundamental quantity: the luminosity function.
The gamma-ray luminosity function, $\Phi_\gamma(z, L_\gamma)$, gives the comoving number density of sources per unit luminosity $L_\gamma$ as a function of redshift $z$.
The connection between the two is given by \cite{Fornasa:2015qua}
\aure{this equation is correct, I double checked}
\begin{equation}
	\frac{dN}{dS} = \int_{z_\mathrm{min}}^{z_\mathrm{max}} dz \int_{\Gamma_\mathrm{min}}^{\Gamma_\mathrm{max}} d\Gamma \; \frac{dN}{dV \, dL_\gamma \, d\Gamma}(z, L_\gamma, \Gamma) \; \frac{dV}{dz \, d\Omega} \; \frac{dL_\gamma}{dS},
	\label{eq:dnds_lf}
\end{equation}
where $\Gamma$ is a parameter that characterizes the energy spectrum shape of the source, $dV/(dz \, d\Omega)$ is the comoving volume element per unit redshift and solid angle, and $dL_\gamma/dS$ is the Jacobian relating luminosity to observed flux through the luminosity distance $D_L(z)$.
The integrand $dN/(dV \, dL_\gamma \, d\Gamma)$ --- the comoving density of sources per unit luminosity and spectral index --- is the luminosity function itself, which encodes the intrinsic abundance of sources at each cosmic epoch.

This integral makes precise the physical content of $dN/dS$: it takes the intrinsic distribution of sources in luminosity and redshift, projects it through the flux--distance relation, and integrates over all redshifts and spectral shapes to produce the observed flux distribution.
In other words, $dN/dS$ can be interpreted as the integral along the line of sight of the gamma-ray luminosity function. It is what we observe when we look at the sky, while the luminosity function is the astrophysical quantity we ultimately want to constrain.

The source-count distribution therefore bridges two distinct observational regimes.
At the bright end, it is anchored by direct measurements from the Fermi-LAT catalog, which shows a flux distribution broadly compatible with an $S^{-2}$ power law (shallower than the Euclidean expectation $S^{-5/2}$) across most of the resolved range \cite{Amerio:2023uet}.
At the faint end, below the detection threshold of approximately $S_\mathrm{th} \sim 2 \times 10^{-10}$~cm$^{-2}$~s$^{-1}$ in the 1--10 GeV band, the $dN/dS$ must be reconstructed statistically from the collective imprint of unresolved sources on the photon-count map.
The techniques for performing this reconstruction --- and for extending the $dN/dS$ well below the Fermi-LAT threshold --- are the subject of the following subsections.

\aure{maybe add a nothe on that $S^{-5/2}$, where it comes from and why, it can be a footnote as that's a question that pops up often in talks: why do we get $S^{-2}$ and do we expect that?}

\subsection{The 1-Point Photon-Count Distribution}
\label{sec:6.2.2}

Recovering the source-count distribution below the detection threshold requires extracting population information from the collective statistical imprint of unresolved sources on the photon-count map.
The photon-count distribution --- also called the 1-point probability distribution function (1pPDF) --- provides a well-established method for doing so.

The idea is straightforward: given a pixelated sky map, one constructs the histogram of the number of pixels containing exactly $k$ photons.
Different source populations leave distinct imprints on this histogram \cite{Fornasa:2015qua}.
Bright but rare sources, such as blazars just below the detection limit, produce multiple photons during the observation period; because these photons are spatially correlated, they populate pixels with moderately large photon counts, generating non-Poissonian tails in the distribution.
By contrast, faint but numerous sources, such as star-forming galaxies, contribute on average far less than one photon each; their combined signal is statistically indistinguishable from a smooth Poissonian background \cite{Malyshev:2011zi}.
The departures from Poisson statistics therefore encode information about the sub-threshold source population and its $dN/dS$.
This connection is made rigorous through a generating-function formalism \cite{Malyshev:2011zi, Fornasa:2015qua} in which the total photon-count probability factorises into independent contributions from point sources, Galactic foreground, and isotropic background, and the parameters of $dN/dS$ are obtained by fitting the predicted pixel histogram to data via maximum likelihood.

Malyshev and Hogg~\cite{Malyshev:2011zi} pioneered this approach for gamma rays, applying it to 11 months of Fermi-LAT data and constraining the contribution of unresolved AGN-like sources to approximately $17 \pm 2\%$ of the total flux at high Galactic latitudes above 1~GeV.
Zechlin, Cuoco et al.~\cite{Zechlin:2015wdz} improved the method by incorporating the spatial variation of the Galactic foreground and exposure at the pixel level, measuring the $dN/dS$ in the 1--10~GeV band down to fluxes approximately one order of magnitude below the detection threshold.
Subsequent extensions by the same group~\cite{Zechlin:2016pme} and by Lisanti et al.~\cite{Lisanti:2016jub} generalised the technique to multiple independent energy bands spanning 1--171~GeV.
The same pixel-statistics formalism was also adapted for the Non-Poissonian Template Fit (NPTF) applied to the Galactic Center Excess (Section~\ref{sec:4.4}), where it served to distinguish a smooth dark matter signal from a population of unresolved point sources such as millisecond pulsars --- illustrating the broad applicability of photon-count methods across very different science cases.

\subsection{Constraining the Composition of the Unresolved Sky}
\label{sec:6.2.3}

The photon-count methods described above --- the 1pPDF and, as we shall see, the CNN-based approach developed in this chapter --- reconstruct the \emph{total} source-count distribution: the aggregate $dN/dS$ of all point sources, regardless of their astrophysical class.
They measure how many sources exist at each flux level, but they cannot, by themselves, separate the blazar contribution from that of star-forming galaxies, misaligned AGNs, or any other population.
This is a fundamental feature of photon-count statistics: they respond to photons, not to source identities.

A complementary body of work approaches the problem from the opposite direction, attempting to estimate the $dN/dS$ of individual source classes by leveraging multi-wavelength information.
For blazars, the most numerous resolved population, the gamma-ray luminosity function can be constructed directly from Fermi-LAT catalog sources by fitting the observed flux distribution with a broken power law and correcting for detection efficiency \cite{Ajello:2015mfa}.
These analyses indicate that unresolved blazars could account for roughly 20--30\% of the UGRB intensity below 100~GeV, while high-synchrotron-peaked BL Lacs dominate the background above a few tens of GeV \cite{Ajello:2015mfa}.
For star-forming galaxies, which are too faint to be resolved in significant numbers, the strategy relies on empirical scaling relations between gamma-ray and infrared luminosity \cite{Fermi-LAT:2012nqz}; the well-measured infrared luminosity function is then used to extrapolate the gamma-ray contribution, yielding estimates that range from a few percent to $\sim\!50\%$ of the UGRB below 10~GeV, depending on the assumed redshift evolution \cite{Fornasa:2015qua, Fermi-LAT:2012nqz}.
Misaligned AGNs (radio galaxies) are treated similarly, through correlations between their core radio luminosity and gamma-ray luminosity \cite{Inoue:2011bm, DiMauro:2014wha}; the resulting estimates place their contribution at roughly 10--25\% of the UGRB, but with an uncertainty band spanning almost an order of magnitude, reflecting the scarcity of the underlying radio data \cite{Fornasa:2015qua}.
Millisecond pulsars contribute less than $1\%$ and are subdominant at all energies.
These numbers should be taken as order-of-magnitude indications rather than precise measurements: the scaling relations are derived from small samples of nearby, bright objects and extrapolated to high redshifts, and the resulting uncertainty propagates directly into the inferred $dN/dS$ for each class.

Since these known populations appear to collectively account for a large fraction of the UGRB intensity \cite{Ajello:2015mfa, Fornasa:2015qua}, a more precise measurement of the total $dN/dS$ would help constrain the room remaining for additional contributors such as dark matter annihilation.

Any such measurement, however, is bounded from below by the detector exposure.
Below a characteristic flux $F_\mathrm{sens} = 1 / \mathcal{E}_\mathrm{av}$ --- where $\mathcal{E}_\mathrm{av}$ is the average exposure in the region of interest --- a point source contributes fewer than one photon on average, and its signal becomes statistically indistinguishable from smooth Poisson noise.
For 14 years of Fermi-LAT data at high Galactic latitudes in the 1--10~GeV band, this floor lies at approximately $4 \times 10^{-12}$~cm$^{-2}$~s$^{-1}$.
Both the 1pPDF \cite{Zechlin:2015wdz, Zechlin:2016pme} and the CNN-based method developed in this chapter can probe the $dN/dS$ down to $\sim\!5 \times 10^{-12}$~cm$^{-2}$~s$^{-1}$, though the reconstruction uncertainties grow rapidly as this limit is approached.

Beyond the direct characterization of the unresolved sky, the $dN/dS$ measured in this chapter provides the empirical prior for the probabilistic cataloging framework developed in Chapter~\ref{ch:7}, where it determines the expected flux distribution of sources just below the detection threshold.
The complementary approach of searching for collective dark matter emission through angular cross-correlations with tracers of large-scale structure is addressed in Chapter~\ref{ch:8}.

The 1pPDF approach reviewed in the preceding subsection has been successful, but it operates under assumptions that limit its scope.
The generating-function framework requires a parametric model for the $dN/dS$ (typically a multiply broken power law), and the inference is not amortized: each new dataset or parameter configuration requires a full re-convergence of the likelihood.
The next section introduces a simulation-based alternative that relaxes these assumptions and enables non-parametric, amortized inference of the source-count distribution.
